\begingroup
\makeatletter\tagsleft@true\let\veqno\@@leqno\makeatother
\newcommand{\TransBracket}[1]{\boldsymbol{[}#1\boldsymbol{]}}
\newcommand{\TransCrescent}{\tikz[baseline=-.1em,x=.09em,y=.09em]{\draw[line width=.4pt] (5,0) .. controls (-1,0) and (-1,8) .. (5,8) .. controls (0,7.5) and (0,.5) .. (5,0);}}
\ArithmeticHead{292}{L}{ANALYSIS. NACHLASS.}
\NumberedItem{4}
\noindent{\small[S. 43]}\par
\par\indent Wir bezeichnen
\[(1-x)(1-xx)(1-x^3)\ \text{etc.}\ \text{mit}\ \TransBracket{x}[\text{*)}],\]
\noindent so ist
\[x\TransBracket{x^{24}}=x-x^{25}-x^{49}+x^{121}+x^{169}-x^{289}-x^{361}-\text{etc.}[\text{**)}],\]
\begin{equation*}\tag*{$[2]'$}
\left\{\begin{aligned}
&1+x+x^3+x^6+\text{etc.}\\
&=\frac{(1-xx)(1-x^4)(1-x^6){.}{.}{.}}{(1-x)(1-x^3)(1-x^5){.}{.}{.}}=\frac{\TransBracket{xx}^2}{\TransBracket{x}}\\
&=(1+x)(1+xx)(1+x^3)(1+x^4)(1+x^5){.}{.}{.}\\
&\qquad.(1-x^2)(1-x^4)(1-x^{16}){.}{.}{.}
\end{aligned}\right.
\end{equation*}
\begin{equation*}\tag*{[3]}
\begin{aligned}
&1-2x+2x^4-2x^9+\text{etc.}\\
&=\frac{(1-x)(1-xx)(1-x^3){.}{.}{.}}{(1+x)(1+xx)(1+x^3){.}{.}{.}}=\frac{\TransBracket{x}^2}{\TransBracket{xx}}
\end{aligned}
\end{equation*}
\begin{equation*}\tag*{[4]}
\left\{\begin{aligned}
&1+2x+2x^4+2x^9+\text{etc.}\\
&=\frac{(1+x)(1-xx)(1+x^3){.}{.}{.}}{(1-x)(1+xx)(1-x^3){.}{.}{.}}=\frac{\TransBracket{xx}^5}{\TransBracket{x}^2\TransBracket{x^4}^2}\\
&=(1+x)^2(1-xx)(1+x^3)^2(1-x^4)(1+x^5)^2{.}{.}{.}[\text{***)}]
\end{aligned}\right.
\end{equation*}
\begin{equation*}\tag*{[5]}
\TransBracket{-x}=\frac{\TransBracket{xx}}{\TransBracket{x}\TransBracket{x^4}}.
\end{equation*}
\ClosingRule
\Needspace{5\baselineskip}
\par\smallskip\noindent\rule{.25\linewidth}{.4pt}\par
{\small
\par\indent [*) Zum Unterschied von den die Zusätze des Herausgebers einschließenden Klammern sind die hier von \textsc{Gauss} benutzten Klammern verstärkt.]
\par\indent [**)] Folgt aus \# [S. 293], wenn man statt $x,y$ schreibt $x^3,-x$ [und dann an die Stelle von $x$ setzt $x^{12}$].
\par\indent [***)] Entwickelt die Potenzen der Reihe.
\par}
\clearpage
\ArithmeticHead{293}{R}{{\scriptsize ZUR THEORIE DER TRANSSCENDENTEN FUNCTIONEN. I}}
\SourceLeaf{300}
\[\begin{aligned}
&{.}{.}{.}\left[\left(1+\frac{x^{16}}{yy}\right)\left(1+\frac{x^8}{yy}\right)\right]\left(y+\frac1y\right)\\
&\qquad(1+x^8yy)(1+x^{16}yy)(1+x^{24}yy){.}{.}{.}\\
&=\frac1{\TransBracket{x^8}x}\left\{\left(y+\frac1y\right)x+\left(y^3+\frac1{y^3}\right)x^9+{.}{.}{.}\right\}
\end{aligned}\]
\[\begin{aligned}
&(1+xy)(1+x^2y)(1+x^3y)(1+x^4y){.}{.}{.}\\
&\qquad\left(1+\frac xy\right)\left(1+\frac{x^2}y\right)\left(1+\frac{x^3}y\right)\left(1+\frac{x^4}y\right){.}{.}{.}\\
&=\frac1{\TransBracket{x}}\Biggl\{1+x\left(y-1+\frac1y\right)+x^3\left(yy-y+1-\frac1y+\frac1{yy}\right)\\
&\qquad+x^6\left(y^3-yy+y-1+\frac1y-\frac1{yy}+\frac1{y^3}\right){.}{.}{.}\Biggr\}
\end{aligned}\]
\ClosingRule
\[\TransBracket{x^3}=1-3x+5x^3-7x^6+9x^{10}-\text{etc.}\]
\noindent folgt leicht aus \TransCrescent\ [S. 294], wenn man $y=1+\omega$ setzt und daraus die Bedingungsgl[eichung] bildet.
\ClosingRule
\begin{equation*}\tag*{\#}
\left\{\begin{aligned}
&(1+xy)(1+x^3y)(1+x^5y){.}{.}{.}\left(1+\frac xy\right)\left(1+\frac{x^3}y\right)\left(1+\frac{x^5}y\right){.}{.}{.}\\
&=\frac1{\TransBracket{xx}}\left\{1+x\left(y+\frac1y\right)+x^4\left(yy+\frac1{yy}\right)+x^9\left(y^3+\frac1{y^3}\right)+{.}{.}{.}\right\},
\end{aligned}\right.
\end{equation*}
\[\begin{aligned}
&\text{etc.}+x^{(\omega-1)^2}+x^{\omega\omega}+x^{(\omega+1)^2}+\text{etc.}\\
&=\frac{\TransBracket{xx}{.}{.}{.}1+x^{2\omega+3}.1+x^{2\omega+1}.1+x^{2\omega-1}.1+x^{2\omega-3}{.}{.}{.}}{x^{-\omega\omega}.x^{2\omega-1}.x^{2\omega-3}{.}{.}{.}}\mathpunct{\raisebox{.5ex}{.}}
\end{aligned}\]
\clearpage
\endgroup
